%! ~~~ Packages Setup ~~~ 
\documentclass[]{../../../../tex/classes/styledArticle}
\usepackage{../../../../tex/packages/hebrewSupport}
\usepackage{../../../../tex/packages/mathShortcuts}
\usepackage{../../../../tex/packages/theoremsSupport}
\usepackage{../../../../tex/packages/enfancysections}

%! ~~~ Document ~~~

\author{שחר פרץ}
\title{\textit{לינארית 2א 23}}
\date{30 ביוני 2025}
\begin{document}
	\maketitle
	\textbf{תרגיל. }חשבו: (כולל פתרון)
	\[ \pms{-1 & 0 \\ 0 & -1}\pms{a & b \\ -b & a}\pms{1 & 0 \\ 0 & -1} = \pms{a & -b \\ b & a} \]
	מכאן נסיק שאכן המטריצות להלן דומות עד לכדי שינוי בסיס, וזו הסיבה שלא איכפת לנו מהסימן של $b$ כמו שראינו בהרצאה 22. 
	
	\theo{תהי $T \co V \to V$ צמודה לעצמה ואי שלילית $\mut{Tv}{v} \ge 0$, אז קיימת ויחידה $R \co V \to V$ אי־שלילית צמודה לעצמה כך ש־$R^2 = T$. }
	
	\begin{proof}
		\textbf{קיום. }
		מהמשפט הספקטרלי קיים בסיס א''נ של ו''ע להעתקה אי־שלילית כל הע''ע הם אי־שליליים. 
		\[ [T]^{B}_B = \diag(\lg_1 \dots \lg_n) \quad [R]_B^B = \diag(\sqrt \lg_1 \dots \sqrt \lg_n) \]
		(ראינו זאת בתרגול). עוד נבחין ש־$R$ צמודה לעצמה כי ע''ע ממשיים. 
		
		\textbf{יחידות. }נבחין שכל ו''ע של $T$ הוא ו''ע של $R$: יהי $i \in [n]$, ו־$B = (e_1 \dots e_n)$ בסיס מלכסן, ואז עבור $R$ צמודה לעצמה כלשהי מתקיים: אז ו''ע של $R$ עם ע''ע $\sqrt \lg$ הוא ו''ע של $T$ עם ע''ע $\lg$ כי: 
		\[ \lg v = R^2v = Tv \implies Rv = \sqrt \lg \]
		הגרירה נכונה מאי־שליליות $R$ שהמשפט מניח עליה יחידות. כלומר הערכים העצמיים של $R$ כלשהי (לא בהכרח זו שברחנו בהוכחת הקיום) נקבעים ביחידות מע''ע של $T$. בסיס של ו''ע של $T$ הוא בסיס ו''ע של $R$, סה''כ ראינו איך $R$ פועלת על בסיס ו''ע כלשהו של $T$ מה שקובע ביחידות את $R$. 
	\end{proof}
	
	\noti{את ה־$R$ לעיל נסמן $\sqrt{T} := R$. }
	
	\section{\en{Polar decomposition}}
	\theo{(פירוק הפולארי) תהי $T \co V \to V$ הפיכה, אז קיימות $R \to V \to V$ חיובית וצמודה לעצמה ו־$U \co V \to V$ אוניטרית כך ש־$T = RU$. }
	\textit{הערה: }לא הנחנו $T$ צמודה לעצמה. הפירוק נכון להעתקה הפיכה כללית. 
	\begin{proof}
		נגדיר $S = TT^*$. נבחין ש־$S$ צמודה לעצמה וחיובית: 
		\[ \forall V \ni v \neq 0 \co \mut{Sv}{v} = \mut{TT^* v}{v} = \mut{T^* v}{T^*v} = \norm{T^*v} > 0 \]
		האי־שוויון האחרון נכון כי $\ker T = \{0\}$, ממשפט קודם $\ker T^* = \ker T = \{0\}$, ו־$v \neq 0$. יצא שזה חיובי ולכן בפרט ממשי, כלומר היא צמודה לעצמה וחיובית. 
		
		קיימת ויחידה $R \co V \to V$ צמודה וחיובית כך ש־$S = R^2$. כל ערכיה העצמיים של $R = \sqrt S$ אינם $0$, ולכן היא הפיכה (ראינו בהוכחה של קיומה שהיא לכסינה יחדיו עם $S$). 
		
		נגדיר $U = R\op T$. נותר להראות ש־$U$ אוניטרית. 
		\[ U^*U = (R\op T)^*(R\op T) = T^*\underbrace{(R\op)^*}_{R\op}R\op T = T^*(R\op)^2 T = T^*S\op T = T^*(TT^*)\op T = I \]
		כדרוש. הטענה $(R\op)^* = R\op$ נכונה משום ש־$R$ צמודה לעצמה. 
	\end{proof}
	
	\textit{הערה לגבי יחידות. }אם $T$ אינה הפיכה, מקבלי םש־$R$ יחידה אבל $U$ אינה. בשביל לא הפיכות נצטרך להצטמצם לבסיס של התמונה ועליו לפרק כמתואר לעיל. במקרה של הפיכות אז $T = RU = R \tl U$ ואז נקבל $R$ הפיכה כלומר $U = \tl U$ וגם $U$ הפיכה. 
	
	עתה נראה ש־$R$ נקבעת ביחידות (בניגוד ליחידות $U$ – יחידות $R$ נכונה גם בעבור פירוק פולארי של העתקה שאיננה הפיכה): 
	\begin{proof}
		\[ TT^* = RU(RU)^* = RUU^*R^* = R^2 \]
		כלומר $R$ היא בכל פירוק שורש, והראינו קודם את יחידות השורש. 
	\end{proof}
	\textit{הערה. }קיים גם פירוק כנ''ל מהצורה $T = UR$. 
	\begin{proof}
		באותו האופן שפירקנו את $T$, נוכל לפרק את $T^* = \tl R \tl U$ פירוק פולארי. נפעיל $^*$ על שני האגפים ונקבל: 
		\[ T^* =\tl R \tl U \implies T = (T^*)^* = \tl U^*\tl R^* = \tl U\op \tl R \]
		נסמן $\tl R =: R, \ \tl U\op =: U$ וסה''כ $T = UR$ כדרוש. 
	\end{proof}
	\lem{עבור $T \co V \to V$ אז ל־$TT^*, T^*T$
		נגדיר $S = TT^*$. נבחין ש־$S$ צמודה לעצמה וחיובית: 
		\[ \forall V \ni v \neq 0 \co \mut{Sv}{v} = \mut{TT^* v}{v} = \mut{T^* v}{T^*v} = \norm{T^*v} > 0 \] יש אותם הערכים העצמיים. }
	\begin{proof}ניעזר בפירוק הפולארי: 
		\begin{align*}
			TT^* &= RUU^*R^* \\
			&= R^2 \\
			TT^* &= U\op R^2U
		\end{align*}
		סה''כ $TT^*, T^*T$ הן העתקות דומות ולכן יש להן את אותם הערכים העצמיים. 
	\end{proof}
	
	\textit{הערה. }אז איך זה קשור לפולארי? $R$ האי־שלילית היא ``הגודל'', בעוד $U$ האוניטרית לא משנה גודל – היא ה''זווית''. 
	
	
	\subsection{פירוק פולארי בעבור מטריצות}
	\theo{(פירוק פולארי עבור מטריצות) תהי $A \in M_n(\F)$ הפיכה, אז קיימות $U, R \in M_n(\F)$ כאשר $U$ א''נ ו־$R$ חיובית צמודה לעצמה כך ש־$A = UR$. }\begin{proof}
		נסתכל על $A^*A$. היא חיובית וצמודה לעצמה (בדומה לעיל). אז $A^*A = P\op D P$, כאשר $D$ אלגסונית חיובית. כאשר $R = P\op \sqrt D P, \ R^2 = AA^*$. היא קיימת ויחידה מאותה הוכחה בדיוק להעתקות. 
	\end{proof}
	
	\section{\en{Singular Value Decomposition (SVD)}}
	\theo{(פירוק לערכים סינגולריים למטריצה – SVD) לכל מטריצה $A \in M_n(\F)$ קיימות מטריצות אוניטריות $U, V$ ומטריצה אלכסונית עם ערכים אי־שלילייים כך ש־$A = UDV$. }
	\begin{proof}
		ידוע שניתן לכתוב $A = \tl UR$ פירוק פולארי. משום ש־$R$ צמודה לעצמה, ניתן לפרקה ספקטרלית ל־$V$ אוניטרית ו־$D$ אלכסונית אי־שלילית (כי $R$ אי־שלילית) כך ש־$R = V\op DV$. סה''כ: 
		\[ A = \underbrace{\tl UV\op}_{=: U} DV = UDV \quad \top \]
		כי $\tl UV\op$ מכפלה של אוניטריות ולכן $U$ אוניטרית כנדרש. 
	\end{proof}
	
	\textit{הערה. }
	\begin{align*}
		AA^* &= (UDV) V^*D^*U^* = UD^2U\op \\
		A^*A &= V\op D^2 V
	\end{align*}
	
	\defi{הערכים העצמיים האי־שליליים של $A^*A$ נקראים \textit{הערכים הסינגולריים} והם נקבעים ביחידות ע''י $A$. }
	הערכים הסינגולרים הם גם הע''ע של $R^2$ הפירוק בפולארי וכן הע''ע של $D^2$ בפירוק SVD. 
	
	\textit{הערה. }פירוק SVD יחיד למטריצה הפיכה. 
	
	{\dotfill \\ \vfil {\begin{center}
				{\Large \textbf{\textit{שחר פרץ $\sim$ סוף הקורס $\sim$ 2025}} \\
					\normalsize הקובץ לא נגמר – יש הרחבה על דואלים בעמוד הבא
					\\
					\scriptsize \textit{קומפל ב־}\en{\LaTeX}\,\textit{ ונוצר באמצעות תוכנה חופשית בלבד}}
		\end{center}} \vfil	}
	\pagebreak
	\defi{בהינתן $V$ מ''ו מעל $\F$, נגדיר $V^* = \hom(V, F)$. }
	\textbf{הבנה. }אם $\dim V = n$ אז $\dim V^* = n$. לכן $V \cong V^*$. לא נכון במקרה הסוף ממדי. 
	
	\lem{יהי $B = (v_i)_{i = 1}^{n}$ בסיס ל־$V$. אז \hfill $\forall i \in [n] \co \exists \psi_i \in V^* \co \forall j \in [n] \co \psi_i(v_j) = \dg_{ij}$}
	
	\theo{יהי $V$ נ''ס ו־$B = (v_i)_{i = 1}^{n}$ אז קיים ויחיד בסיס $B^* = (\psi_i)_{i = 1}^{n}$ המקיים $\forall i, j \in [n] \co \psi_i(v_j) = \dg_{ij}$. }\begin{proof}
		נבחין שהבדרנו העתקה לינארית $\phi \co B \to V^*$ והיא מגדירה ביחידות $\psi$ לינארית   $\psi \co V \to V^*$ המקיימת את הנרש. ברור שהבנייה של $\phi_i$ קיימת ויחידה כי היא מוגדר לפי מה קורה לבסיס. נותר להוכיח שזה אכן בסיס. יהיו $\ag_1 \dots \ag_n \in \F$ כך ש־$\sum\ag_i \psi_i = 0$. (האפס הזה הוא פונקציונל האפס). יהי $j \in [n]$. אז $0(v_j) = 0 = \cl{\sumni \ag_i \psi_i}v_j = \sumni \ag_i \psi_i(v_j) = \sum\ag_i \dg_{ij} = \ag_j$ וסה''כ $\ag_j = 0$. 
	\end{proof}
	
	נבחין שאפשר להגדיר: 
	\defi{$V^{**} = \hom(V^*, \F)$}
	ואכן $\dim V < \inf$ אז: 
	\[ V \cong V^* \cong V^{**} \]
	במקרה הזה, בניגוד לאיזו' הקודם, יש איזו' ``טבעי'' (קאנוני), כלומר לא תלוי באף בסיס. 
	
	\theo{קיים איזומורפיזם קאנוני בין $V$ ל־$V^{**}$. }
	\begin{proof}נגדיר את האיזו' הבא: 
		\[ \psi \co V \to V^{**} \quad \psi(v) =: \bar v \quad \forall \psi \in V^* \co \bar v(\psi) = \psi(v) \]
		נוכיח שהוא איזו': 
		\begin{itemize}
			\item \textbf{ט''ל: }יהיו $\ag, \bg \in \F, \ v, u \in V$. אז: 
			\[ \psi(\ag v + \bg u) = \overline{\ag v + \bg u} \seq \ag \bar v + \bg \bar u \]
			נוכיח זאת: 
			\[ \overline{\ag v + \bg u}(\phi) = \phi(\ag v + \bg u) = \ag \phi(v) + \bg \phi(u) = \ag v(\phi) + \bg \bar u(\phi) = (\ag \bar v + \bg \bar u)(\phi) = (\ag \psi(v) + \bg \psi(u))(\phi) \]
			\item \textbf{חח''ע: }יהי $v \in \ker \psi$. רוצים להראות $v = 0$. 
			\[ \forall \phi \in V^* \co \bar (\phi) = 0 \implies \forall \phi \in V^* \co \phi(v) = 0 \]
			אם $v$ אינו וקטור האפס, נשלימו לבסיס $V = (v_i)_{i = 1}^{n}$ ואם $\phi_1 \dots \phi_n$ בסיס הדואלי אז $\phi_1(v) = 1$ אבל אז $0 = \bar v(\phi_1) = 1$ וסתירה. 
			\item \textbf{על: }משוויון ממדים $\dim V^{**} = \dim V$. 
		\end{itemize}
	\end{proof}
	כלומר, הפונקציונלים בדואלי השני הם למעשה פונקציונלים שלוקחים איזשהו פונקציונל בדואלי הראשון ומציבים בו וקטור קבוע. 
	
	
	\noti{לכל $v \in V$ ו־$\phi \in V^*$ נסמן $\phi(v) = (\phi, v)$}
	\theo{ייהו $V, W$ מ''וים נוצרים סופית מעל $\F$, $T \co V \to W$. אז קיימת ויחידה $T^* \co W^* \to V^*$ כך ש־$(\psi, T(v)) = (T^*(\psi), v)$. }
	אם לצייר דיאגרמה: 
	\[ V \overset{T}{\to} W \cong W^* \overset{T^*}{\to} V^* \cong V \]
	(תנסו לצייר את זה בריבוע, ש־$V, W$ למעלה ו־$V^*, W^*$ למטה, כדי להבין ויזולאית למה זה הופך את החצים)
	
	ברמה המטא־מתמטית, בתורת הקטגוריות, יש דבר שנקרא פנקטור – דרך לזהות בין אובייקטים שונים במתמטיקה. מה שהוא עושה, לדוגמה, זה להעביר את $\hom(V, W)$ – מרחבים וקטרים סוף ממדיים – למרחב המטריצות. הדבר הזה נקרא פנקטור קו־וראיינטי. במקרה לעיל, זהו פנקטור קונטרא־ווריאנטי – שימוש ב־$T^*$ הופך את החצים. (הרחבה של המרצה)
	
	אז אפשר להגדיר פנקטור אבל במקום זה נעשה את זה בשפה שאנחנו מכירים – לינארית 1א. בהינתן $\psi \in W^*$, נרצה למצוא $T^*(\psi) \in V^*$. נגדיר: 
	\[ T^*(\psi) = \psi \circ T \]
	ברור מדוע $T^* \co W^* \to V^*$. בעצם, זהו איזומורפיזם (``בשפת הפנקטורים'') קאנוני. עוד קודם לכן ידענו (בגלל ממדים) שהם איזומורפים, אך לא מצאנו את האיזומורפיזם ולא ראינו שהוא קאנוני. 
	\[ \tau \co \hom(V, W) \to \hom(W^*, V^*) \quad \tau(T) = T^*  \]
	היא איזומורפיזם. 
	
	(הערה: תודה למרצה שנענה לבקשתי ולא השתמש ב\slash phi אחרי שעשיתי \slash renewcommand \slash phi \{\slash varphi\})
	\begin{proof}[הוכחת לינאריות]
		יהיו $T, S \in \hom(V, W)$, $\ag \in \F$. אז: 
		\[ \tau(\ag T + S) = (T + \ag S)^* \]
		יהי $\psi \in W^*$, אז: 
		\[ (T + \ag S)^*(\psi) = \psi \circ (T + \ag S) \]
		יש למעלה פונקציונל ב־$V^*$. ננסה להבין מה הוא עושה על $V$. יהי $v \in V$: 
		\[ [\psi (T + \ag S)](v) = \psi((T + \ag S)v) = \psi(T)v + \ag \psi(S)(v) = (\psi\circ T + \ag \psi\circ S)(v) = ((T^* + \ag S^*)\circ (\psi))v = (\tau(T) + \ag \tau(S))(\psi)(v) \]
		סה''כ קיבלנו לינאריות: 
		\[ \tau(T + \ag S) = \tau(T) + \ag \tau(S) \]
	\end{proof}
	נוכל להוכיח זאת יותר בפשטות עם הנוטציה של ``המכפלה הפנימית'' שהגדרנו לעיל, $(\phi, v)$. 
	עתה נוכיח ש־$\tau$ לא רק לינארית, אלא מוגדרת היטב. 
	\begin{proof}\,
		\begin{itemize}
			\item \textbf{חח''ע: }תהי $T \in \ker \tau$, אז $\tau(T) = T^* = 0_{\hom(W^*, V^*)}$. נרצה להראות ש־$T$ העתקה האפס. נניח בשלילה ש־$T \neq 0$. אז קיים $v' \in V$ כך ש־$T(v') \neq 0$. נשלימו לבסיס – $(T(v) = w_1, w_2 \dots w_n)$ בסיס ל־$W$. יהי $(\psi_1 \dots \psi_n)$ הבסיס הדואלי. אז:
			\[ \tau(T)(\psi_1) = T^*(\psi_1) = \psi_1  \]
			אז: 
			\[ 0 = \tau(T)(\psi_1)(v') = T^*(\psi_1)(v') = \psi_1 \circ T(v') = \psi_1(w_1) = 1 \]
			סתירה. לכן $\ker \tau = \{0\}$ ולכן $\tau$ חח''ע. 
			\item \textbf{על: }גם כאן משוויון ממדים
		\end{itemize}
	\end{proof}
	
	\textbf{שאלה ממבחן שבן עשה. }(``את השאלה הזו לא פתרתי בזמן המבחן, ואני די מתבייש כי אפשר לפתור אותה באמצעות כלים הרבה יותר פשוטים'' ``חה חה'' ``לא חח''ע זה חד־חד ערכי'') יהיו $V, W$ מ''ו מעל $\F$ ו־$(w_1 \dots w_n)$ בסיס של $W$. תהי $T \co V \to V$. הוכיחו שקיימים $\phi_1 \dots \phi_n \in V^*$ כך שלכל $v \in V$ מתקיים: 
	\[ T(v) = \sumni \phi_i(v) w_i \]
	\textbf{שימו לב: }בניגוד למה שבן עשה במבחן, $V$ לא בהכרח נוצר סופית. 
	\begin{proof}[הוכחת ראש בקיר.]
		לכל $v \in V$ קיימים ויחידים $\ag_1 \dots \ag_n$ כך ש־$T(v) = \sumni \ag_i w_i$. נגדיר $\forall i \in [n] \co \phi_i(v) = \ag_i$. זה לינארי. 
	\end{proof}
	
	\begin{proof}[הוכחה ``מתוחכמת''.]
		``אני אהבתי את ההוכחה שלי'': נתבונן בבסיס הדואלי $B^* = (\psi_1 \dots \psi_n)$ שמקיים את הדלתא של קרונקר והכל. 
		נגדיר $T^*(\psi_i) =: \phi_i$. יהי $v \in V$. קיימים ויחידים $\ag_1 \dots \ag_n$ כך ש־$T(v) = \sumni \ag_i w_i$. אז: 
		\[ \sumni \phi_i(v)w_i = \sum T^*(\psi_i)(v) w_i \]
		צ.ל. $\ag_i = T^*(\psi_i)(v)$. אך נבחין שהגדרנו: 
		\[ T^*(\psi_i)(v) = \psi_i(T(v)) = \psi_i\cl{\sumni \ag_j w_j} = \ag_j \]
		
	\end{proof}
	
	``הפכת למרצה במתמטיקה כדי להתנקם באחותך?'' ``כן.''
	
	\subsection{המאפס הדואלי ומרחב אורתוגונלי}
	\defi{יהי $V$ מ''ו נוצר סופית. יהי $S \subseteq V$ קבוצה. נגדיר $\{\phi \in V^* \mid \forall v \in S \co \phi(v) = 0\} =: S^0 \subseteq V^*$. }
	
	\textbf{דוגמאות. } \hfil $\{0\}^0 = V^*, \ V^0 = \{0\}$
	
	\theo{\,
	\begin{enumerate}
			\item $S^0$ תמ''ו של $V^*$. 
			\item \hfil $(\Sp S)^0 = S^0$
			\item \hfil $S \subseteq S' \implies S'^0 \subseteq S^0$
	\end{enumerate}}

	\theo{יהי $V$ נ''ס, $U \subseteq V$ תמ''ו. אז $\dim U + \dim U^0 = n$}
	באופן דומה אפשר להמשיך ולעשות: 
	\[ \dim U^0 + \dim U^{**} = n \]
	וכן: 
	\[ U \cong U^{**} \]
	איזומורפיזם קאנוני. זאת כי:
	\[ \forall \phi \in U^0 \, \forall u \in U \co \phi(u) = 0 \]
	ומי אלו הוקטורים שיאפסו את $\phi$ שמאפס את $u$? הוקטורים ב־$U$ עד לכדי האיזומורפיזם הקאנוני מ־$U$ ל־$U^{**}$. 
	
	נבחין ש־: 
	\[ [T^*]^{\bc^*}_{\ac^*} = ([T]^{\ac}_{\bc})^T \]
	כאשר $\ac$ בסיס ל־$V$, $\ac^*$ ל־$V^*$, $\bc$ ל־$W$, $\bc^*$ ל־$W^*$. 
	
	``כוס אמא של קושי'' – בן על זה שקושי גילה את המשפט לפניו. 
	
	
	\ndoc
\end{document}
